%==============================================================================
\section{Introduction}
\label{sec:intro}
%==============================================================================

% P1: Problem — repo-level file retrieval, importance of negatives
Resolving real-world software issues in large repositories requires accurately localizing the relevant source files from hundreds or thousands of candidates~\cite{jimenez2024swebench}.
Dense retrieval has emerged as the dominant paradigm for this task, with recent systems employing contrastive learning to train encoders that map issue descriptions and code files into a shared embedding space~\cite{reddy2025swerank,xie2025swefixer}.
The quality of these retrievers depends critically on the selection of \emph{negative} examples during training: which non-relevant files the model learns to distinguish from the ground truth.
Hard negative mining---selecting challenging negatives that expose the model to confusable examples---has proven essential for training effective dense retrievers in open-domain text retrieval~\cite{karpukhin2020dpr,xiong2021ance,zhan2021star}.
Yet for code retrieval, systematic understanding of what mining strategies work and why remains limited.

% P2: Prior work contradiction — CoCoHaNeRe says embedding negatives harmful; our question
In the code domain, the most relevant prior study is CoCoHaNeRe~\cite{fan2025cocohanere}, which proposes gradient-based hard negative mechanisms for code search and reports that non-differentiable embedding-based negatives are ``slow and ineffective,'' harming retrieval performance.
This finding has influenced subsequent work to avoid embedding-based mining for code.
However, we observe that CoCoHaNeRe used BM25-based retrieval as a proxy for embedding distance---a critical methodological conflation, as BM25 selects lexically similar negatives that may destroy useful token-overlap signals during contrastive training.
This raises a natural question: does \emph{true} embedding-distance mining---using the encoder's own pretrained representations---genuinely fail for code retrieval, or has its potential been obscured by evaluation confounds?

% P3: Code structure as unique dimension + evaluation methodology gap
We identify two underexplored dimensions in this space.
First, code possesses rich structural information---call graphs, data-flow dependencies, inter-procedural taint tracking---absent in natural language text.
This structural connectivity can serve as an independent signal for constructing hard negatives: files that are structurally disconnected from the query yet lexically similar provide training examples grounded in program semantics.
We operationalize this insight using CodeQL~\cite{avgustinov2016ql} inter-procedural data-flow analysis to compute structural connectivity at repository scale.

% P4: Evaluation methodology gap
Second, we identify that \emph{evaluation methodology} critically affects measured strategy effectiveness.
Standard pool-based evaluation---small fixed pools with random negatives, unclear checkpoint selection criteria---can mask genuine strategy effects or create illusory ones.
We find that the choice of checkpoint selection criterion alone (best-by-validation-loss vs.\ best-by-R@1) yields a \placeholder{+16.3pp} R@1 difference, substantially exceeding any strategy effect.
When evaluation controls for these confounds via query-split protocols and realistic pool sizes, a clearer picture of strategy effectiveness emerges.

% P5: Contributions — VERSION-DEPENDENT
% ===========================================================
% VERSION A (criterion dominates):
% Our systematic comparison of four negative mining strategies---random, BM25,
% embedding-distance, and dataflow-aware---across three encoders yields the
% following findings:
% (1) Checkpoint selection criterion is the dominant factor in contrastive
%     retrieval evaluation (+16.3pp), exceeding all strategy effects.
% (2) A query-split evaluation protocol reveals effects masked by standard eval.
% (3) When evaluation is properly controlled, embedding mining shows marginal
%     but consistent benefit.
% ===========================================================
% VERSION B (mining works + methodology):
% Our systematic comparison yields:
% (1) Simple static embedding-distance mining improves code retrieval by
%     \placeholder{X.Xpp} at realistic pool sizes, an effect masked by
%     small-pool evaluation and inappropriate checkpoint selection.
% (2) A query-split evaluation protocol with explicit checkpoint criterion
%     reporting reveals effects obscured by standard methodology.
% (3) Structure-grounded (dataflow) negatives exhibit encoder-dependent
%     interaction patterns mediated by representation geometry.
% ===========================================================

\placeholder{[P5: Version-dependent contribution paragraph. Final version determined after D062 14/14 completion.]}

%==============================================================================
